\section{Problem Formulation}
\label{sec:problem}

\paragraph{Setting.}
A prompt is a pair $(q, C)$: a question $q$ written on the user side and an excerpt $C$
drawn from an IoT log. The excerpt is split into units $U(C)=\{u_1,\dots,u_m\}$, one per
record or key--value group. Each case carries a set of protected values
$V=\{v_1,\dots,v_n\}$; a unit is \emph{sensitive} when it contains a value in $V$. A model
$M$ answers $q$ from the compressed context, and the answer is scored against a reference
span.

\paragraph{What must not reach the model.}
The prompt is the object of study. It travels from the device to the model process and may
be logged, cached, or inspected. A policy is \emph{leak-free} on a case when no value in $V$
is recoverable from the exact text handed to the model, that is, from the rewritten question
$q'$ concatenated with the compressed context, not from the context alone. Recoverability is
measured with the pattern set of the published protocol (full value plus fragments). We
assume no adversarial manipulation of the model and no side channel beyond the prompt text,
and we do not model an adversary who holds the alias map; Section~\ref{sec:discussion}
states what that leaves out.

\paragraph{The queried entity.}
Each question names an entity $e^\star$, which may be an account, a device, or a record
identifier. The \emph{anchor} $A\subseteq U(C)$ is the set of units containing a mention of
$e^\star$. A unit plays a referential role when the question cannot be resolved without
knowing which entity a value belongs to.

\paragraph{Budget.}
A budget arm receives a ratio $r$ and must realise a compressed context of at most
$B=\lceil r\cdot n_C\rceil$ tokens, where $n_C$ is the token length of the uncompressed
excerpt under the tokenizer used to build the prompt. Exceeding $B$ is a hard error rather
than a silent truncation, because in this domain the answer region is often at the end of the
excerpt and truncation would remove it without notice.

\paragraph{Metrics.}
\begin{itemize}
\item \textbf{Task quality.} EM and F1 from the standard SQuAD~v2 scorer;
\textsc{unknown} is not a legal answer.
\item \textbf{Prompt leak}. Fraction of cases in which any protected value is recoverable
from question plus compressed context. Lower is better.
\item \textbf{Gold-in-context (GIC)}. Fraction of cases in which the reference span survives
compression. It is model-free, and therefore an upper bound on what the reader can achieve;
it is also blind to whether the surviving span is still interpretable, which is why we report
task quality as a separate endpoint.
\item \textbf{Overshoot}. Indicator that a budget arm exceeded $B$ on a case.
\end{itemize}

\paragraph{The condition we study.}
The setting is narrow and we state it as such: long, identifier-dense operational text; a
question that names an entity; and a requirement that the entity's value not reach the
model. Questions of this form are common in IoT operations and, as
Section~\ref{sec:analysis} shows, are exactly the ones on which redundancy-based compression
is harmful. The same routing reduces to ordinary compression when the question is topical
and names no entity; we report that variant as out of scope rather than as a success.
